\section{Purpose, scope and the meaning of the comparison}
\label{sec:introduction}

Negative weights in additive NLO matching carry cancellations between
matrix elements, subtraction terms and shower approximations. Removing
them is therefore a change of the matching calculation unless the removed
signed measure is compensated elsewhere. The multiplicative--accumulative
construction of Nason and Salam~\cite{NasonSalam} supplies a useful ideal
limit: use the shower where its leading approximation exceeds the real
matrix element, and retain an additive correction where the real
contribution is larger. The colour-singlet implementation studies of
Refs.~\cite{CheThesis,CheFrederix} motivate inspecting the first-emission
compensation explicitly. Here the object under study is a more general
implementation with signed residuals, ordinary source G/P terms, coloured
Born processes, and an additional Delta source option.

This document describes the branch and completed calculations as they
stand on 24 October 2026. It supersedes neither the frozen earlier paper
nor the cited matching prescriptions. It specifies which realization was
actually used and what its present evidence establishes. In particular,
the name \macnlops{} below refers to this bounded, signed-residual
realization unless an ideal limiting formula is stated explicitly.
The 25 October revision adds the acceptance comparison in
Section~\ref{sec:delta-undamped-audit} and records subsequent computational
optimizations. The numerical results retain the matching prescription
used in the completed campaign.

The conceptual chain is simple but each arrow has a contract:
\begin{equation}
 \begin{gathered}
 \text{qualified common source}
 \longrightarrow \text{first-step proposals}
 \longrightarrow \text{removal/veto}\\
 \longrightarrow \text{accepted LHE}
 \longrightarrow \text{optional continuation}.
 \end{gathered}
 \label{eq:workflow}
\end{equation}
The first arrow uses a shower, the second uses both generated matrix
elements and a leading scalar shower kernel, and the last arrow is a new
shower invocation. They cannot be identified merely because they share a
PDF set or an event filename.

\subsection{Five distinctions that control the interpretation}
First, the negative part of the \emph{generated} hard density is required.
It is not, in general, the negative part of a bare real matrix element
minus a single positive splitting function. The actual source contains
G-function completion, low-scale factors, luminosities and a specific
sampling partition.

Second, the leading positive kernel $K_g$ in the veto denominator is not
the physical Sudakov-suppressed first-emission density. It is also not a
trial overestimate. Its agreement with the leading coefficient of the
native shower is a substantive condition, not a consequence of naming
both objects ``the MC counterterm''.

Third, the first-step plots and full-shower plots have different
operational references. At the first step the curve labelled \mcatnlo{}
is the unrejected matching-driver proposal population, using the same
canonical hard presentation and first-step profile as the accepted
sample. At full shower the black reference is a separate uninterrupted
shower from the original MC-mass LHE body. The blue control restarts the
unrejected first-step endpoints. The comparison thus exposes effects of
the first-step profile and presentation as well as of the restart.

Fourth, an explicit sum over native emitter histories, an explicit sum
over Born-colour flows, and an integration Monte Carlo over outer FKS
sectors are different operations. A sign taken before one of these sums
need not equal the sign taken afterwards. The source qualification is
designed to prevent that distinction from being hidden in a veto ratio.

Fifth, the hard-model and shower masses define different on-shell
phase-space conventions linked by a transport. Keeping the ordinary remassed LHE file as
the reference source does not make its body suitable for an inverse map
derived for the pre-shuffle hard masses. The source annotation preserves
the latter information without rewriting the original event.

\subsection{What is common, and what is independently generated}
Within each process and Delta setting, both methods use the exact same
million-event source. The ordinary matching prescription retains the
ordinary source G terms and bogus factor. The Delta source is a separate
matching prescription with its own compensating contribution. There are
therefore ten independent source cohorts, not one source shared across
ordinary and Delta matching. No claim of event-by-event pairing is made
between those two source settings.

Preparing a source with \code{macnlops=True} adds hard-point/profile
information and selects the supported complete source partition. It is
not a promise to retrofit arbitrary anonymous older LHE files. Nor does
it leave the generator's default stochastic colour/history sampling
unchanged: a global hard-sign reconstruction requires the effective sums
described below. The ordinary \emph{matching formula} is retained; the
qualification and Monte Carlo organization make its signed density
available to the postprocessor.

The implementation uses no MAcNLOPS-specific \pythia{} library patches.
The application interface and the generated density evaluator are
modified. ``Stock'' in this document means the stock-interface runtime
without the retired MAcNLOPS headroom/provider/nf4-coupling extensions; it
does not certify an arbitrary locally installed library as byte-identical
to an upstream release. The recorded library and PDF fingerprints define
the actual calculation. An unrelated local QED modification mentioned in
the development record is outside this algorithm, and QED evolution is
disabled in the validation.

The completed campaign checks finite samples, counts, kinematics,
normalization and the observed distributions. It does not by itself prove
universal positive weights, exact finite-coupling cancellation, a
globally valid native trial envelope, or all-order/logarithmic accuracy.
Those questions must be separated from successful execution of the
declared event-level algorithm.
